\subsection{符号}

14.2.1. 设 \(k\) 是一个 \(\leqslant r\) 的正整数。考虑正合序列

\[
0 \rightarrow \mathrm{P} _ {k} (\mathrm{T} (\mathrm{X}); \mathrm{E}) \xrightarrow {\iota} \mathrm{P} ^ {k} (\mathrm{E}) \xrightarrow {\jmath} \mathrm{P} ^ {k - 1} (\mathrm{E}) \rightarrow 0
\]

（12.6.8），其中我们已令 \(j = r^{k, k-1}\)。把向量函子 M → L(M; F) 应用于这个序列，得到正合序列

\[
0 \rightarrow \mathscr {L} (\mathrm{P} ^ {k - 1} (\mathrm{E}); \mathrm{F}) \xrightarrow {f ^ {\prime}} \mathscr {L} (\mathrm{P} ^ {k} (\mathrm{E}); \mathrm{F}) \xrightarrow {\sigma_ {k}} \mathscr {L} (\mathrm{P} _ {k} (\mathrm{T} (\mathrm{X}); \mathrm{E}); \mathrm{F}) \rightarrow 0
\]

它可以写成：

令

\[
0 \rightarrow \mathrm{D} ^ {k - 1} (\mathrm{E}, \mathrm{F}) \xrightarrow {j ^ {\prime}} \mathrm{D} ^ {k} (\mathrm{E}, \mathrm{F}) \xrightarrow {\sigma_ {k}} \mathrm{S} ^ {k} (\mathrm{E}, \mathrm{F}) \rightarrow 0
\]

\[
\mathbf {S} ^ {k} (\mathrm{E}, \mathrm{F}) = \mathcal {L} (\mathrm{P} _ {k} (\mathrm{T} (\mathrm{X}); \mathrm{E}); \mathrm{F}).
\]

同态 \(j' = \mathcal{L}(j; \mathrm{Id}_{\mathrm{F}})\) 是 \(\mathbf{D}^{k-1}(\mathbf{E}, \mathbf{F})\) 到 \(\mathbf{D}^k(\mathbf{E}, \mathbf{F})\) 的典范包含；同态 \(\sigma_k\) 按定义等于 \(\mathcal{L}(\iota; \mathrm{Id}_{\mathrm{F}})\)。若 D 是一个 \(\mathbf{E} \to \mathbf{F}\) 型、阶 \(\leqslant k\) 的微分算子，则它在 \(\sigma_k\) 下的像是一个截面 \(\sigma_k(\mathbf{D})\)（\(\mathbf{S}^k(\mathbf{E}, \mathbf{F})\) 的截面），我们称之为 D 的 \(k\)-符号（或简称符号）；我们有 \(\sigma_k(\mathbf{D}) = 0\) 当且仅当 D 的阶 \(\leqslant k - 1\)。若 D 是 \(\mathbf{C}^h\) 类的，则其符号也是。

14.2.2（“符号丛的恒同”）。可以把丛 \(S^{k}(E, F) = \mathcal{L}(P_{k}(T(X); E); F)\) 用多种方式写出。首先，若 \(x \in X\)，则 \(P_{k}(T_{x}(X); E_{x})\) 的元素可（13.1.2）等同于 \(\mathcal{L}(\mathsf{T}\mathsf{S}^{k}(\mathsf{T}_{x}(X)); E_{x})\) 的元素。这样我们得到向量丛的恒同

\[
\mathrm{P} _ {k} (\mathrm{T} (\mathrm{X}); \mathrm{E}) = \mathscr {L} (\mathrm{TS} ^ {k} (\mathrm{T} (\mathrm{X})); \mathrm{E}) = (\mathrm{TS} ^ {k} (\mathrm{T} (\mathrm{X}))) ^ {*} \otimes \mathrm{E}
\]

再把向量函子 \(\mathbf{M} \mapsto \mathcal{L}(\mathbf{M}; \mathbf{F}) = \mathbf{M}^* \otimes \mathbf{F}\) 应用于上式，得到：

\[
\mathrm{S} ^ {k} (\mathrm{E}, \mathrm{F}) = \mathrm{T} \mathrm{S} ^ {k} (\mathrm{T} (\mathrm{X})) \otimes \mathrm{E} ^ {*} \otimes \mathrm{F} = \mathrm{T} \mathrm{S} ^ {k} (\mathrm{T} (\mathrm{X})) \otimes \mathscr {L} (\mathrm{E}; \mathrm{F}).
\]

还可以进一步变换上述表达式。若 M 与 N 是两个向量丛，用 \(\operatorname{Sym}^{k}(\mathbf{M};\mathbf{N})\) 表示 \(\mathcal{L}(\mathbf{M},\ldots,\mathbf{M};\mathbf{N})\) 中由对称 \(k\)-线性映射（从 \(\mathbf{M}^{k}\) 到 \(\mathbf{N}\)）组成的子丛。有一个典范同构

\[
\operatorname{Sym} ^ {k} (\mathbf {M}; \mathbf {N}) \rightarrow \operatorname{Sym} ^ {k} (\mathbf {M}; \mathbf {K} _ {\mathbf {X}}) \otimes \mathbf {N};
\]

另一方面，\(\otimes^{k}\mathbf{M}\) 与 \(\otimes^{k}\mathbf{M}^{*}\) 之间的对偶把 \(\operatorname{Sym}^{k}(\mathbf{M};\mathbf{K}_{\mathbf{x}})\) 等同于 \(\operatorname{TS}^{k}(\mathbf{M}^{*})\)。于是得到恒同

\[
\operatorname{Sym} ^ {k} (\mathbf {M}; \mathbf {N}) = \operatorname{TS} ^ {k} (\mathbf {M} ^ {*}) \otimes \mathbf {N}.
\]

把这个公式应用于 M = T(X)* 与 N = ℒ(E; F)，就有：

\[
\mathrm{S} ^ {k} (\mathrm{E}, \mathrm{F}) = \mathrm{T} \mathrm{S} ^ {k} (\mathrm{T} (\mathrm{X})) \otimes \mathscr {L} (\mathrm{E}; \mathrm{F}) = \operatorname{Sym} ^ {k} (\mathrm{T} (\mathrm{X}) ^ {*}; \mathscr {L} (\mathrm{E}; \mathrm{F})).
\]

假设 K 的特征为 0 或 \(>k\)。若 u 是一个对称 \(k\)-线性映射，用 \(\tilde{u}\) 表示多项式映射

\[
x \mapsto \frac {1}{k !} u (x, \dots , x)
\]

这就是相应的多项式映射。映射 \(u \mapsto \tilde{u}\) 定义了一个同构

\[
\mathrm{S} ^ {k} (\mathrm{E}, \mathrm{F}) = \operatorname{Sym} ^ {k} (\mathrm{T} (\mathrm{X}) ^ {*}, \mathscr {L} (\mathrm{E}; \mathrm{F})) \rightarrow \mathrm{P} _ {k} (\mathrm{T} (\mathrm{X}) ^ {*}; \mathscr {L} (\mathrm{E}; \mathrm{F})).
\]

于是，每个元素 \(\sigma\)（属于 \(S^{k}(E,F)_{x}\)，其中 \(x\in X\)）都可以等同于一个齐次多项式映射 \(\tilde{\sigma}\)（\(k\) 次的，从 \(T_{x}(X)^{*}\) 到 \(\mathcal{L}(E_{x};F_{x})\)）。

14.2.3（“微分算子符号的计算”）。设 \(x \in X\)，并设 D 是一个 \(E \to F\) 型、阶 \(\leqslant k\) 的微分算子，它定义在 \(x\) 的一个开邻域中，且是 \(\mathbf{C}^{r-k}\) 类的。我们现在来明确写出值 \(\sigma_k(D)(x)\)（D 的符号在 \(x\) 处的值）。把它看作 \(\text{Sym}^k(T_x(X)^*; \mathcal{L}(E_x; F_x))\) 的元素（参见 14.2.2）。设 \(\omega_1, \ldots, \omega_k\) 是 \(x\) 处的余向量，并选取函数 \(f_1, \ldots, f_k\)（\(\mathbf{C}^r\) 类，定义在 \(x\) 的一个开邻域 U 上），使得 \(d_x f_i = \omega_i\)（对 \(i = 1, \ldots, k\)）。令

\[
\mathbf {D} ^ {\prime} = (- 1) ^ {k} \mathrm{ad} (f _ {k}) \dots \mathrm{ad} (f _ {1}) \mathbf {D} \quad (\text { cf. } 14.1.10, a).
\]

算子 D 的阶 \(\leqslant0\)；它是 \(\mathcal{L}(\mathrm{E};\mathrm{F})|\mathrm{U}\) 的一个截面。\(\mathbf{D}'\) 在 x 处的值是一个元素 \(\lambda(\omega_{1},\ldots,\omega_{k})\)（\(\mathcal{L}(\mathrm{E}_{x};\mathrm{F}_{x})\) 中的元素），它只依赖于 \((\omega_{1},\ldots,\omega_{k})\)。这样定义的映射 \(\lambda\) 是对称的。它等于 D 在 x 处的符号 \(\sigma_{k}(\mathbf{D})(x)\)。

假设 K = R 或 C，并把 \(\sigma_{k}(\mathbf{D})(x)\) 视为从 \(\mathrm{T}_{x}(\mathrm{X})^{*}\) 到 \(\mathcal{L}(\mathrm{E}_{x};\mathrm{F}_{x})\) 的 k 次齐次多项式映射（参见 14.2.2）而把它明确写出。设 \(\omega \in \mathrm{T}_{x}(\mathrm{X})^{*}\)，\(v \in E_{x}\)；选取 x 的开邻域 U 上的一个 \(C^{r}\) 类函数 f，使得 \(d_{x}f = \omega\)，并选取 E 在 U 上的一个 \(C^{r}\) 类截面 s，使得 \(s(x) = v\)。存在 F 在 U 上的截面族 \((\varphi_{0}, \varphi_{1}, \ldots, \varphi_{k})\)，且是唯一的，使得

\[
e ^ {- t f} \mathrm{D} (e ^ {t f} s) = \sum_ {j = 0} ^ {k} t ^ {j} \varphi_ {j} \quad \text { for all } \quad t \in \mathbf {K}.
\]

那么 \(\varphi_{k}(x)\) 与 \(f\) 及 \(s\)（满足 \(d_{x}f = \omega\)、\(v(x) = s\)）的选取无关；映射 \(v \mapsto \varphi_{k}(x)\) 是一个元素 \(\lambda(\omega)\)（\(\mathcal{L}(\mathrm{E}_{x}; \mathrm{F}_{x})\) 中的元素），而 \(\lambda\) 是一个 \(k\) 次齐次多项式映射（从 \(\mathrm{T}_{x}(\mathrm{X})^{*}\) 到 \(\mathcal{L}(\mathrm{E}_{x}; \mathrm{F}_{x})\)）。这个映射等于符号 \(\sigma_{k}(\mathrm{D})(x)\)（D 在 \(x\) 处的符号）。

14.2.4（“标量情形”）。取 \(\mathbf{E} = \mathbf{F} = \mathbf{K}_{\mathbf{x}}\)，这时把 \(\mathbf{D}^k(\mathbf{E}, \mathbf{F})\) 与丛 \(\mathbf{T}^{(k)}(\mathbf{X})\)（阶 \(\leqslant k\) 的点分布的丛）恒同（参见 14.1.6）。我们有

\[
\mathrm{S} ^ {k} (\mathrm{E}, \mathrm{F}) = \mathrm{TS} ^ {k} (\mathrm{T} (\mathrm{X}))
\]

而符号

\[
\sigma_ {k} \colon \mathrm{T} ^ {(k)} (\mathrm{X}) \rightarrow \mathrm{TS} ^ {k} (\mathrm{T} (\mathrm{X}))
\]

不是别的，正是复合

\[
\mathbf {T} ^ {(k)} (\mathbf {X}) \rightarrow \mathbf {T} ^ {(k)} (\mathbf {X}) / \mathbf {T} ^ {(k - 1)} (\mathbf {X}) \xrightarrow {i _ {k}} \mathsf {T S} ^ {k} (\mathbf {T} (\mathbf {X})),
\]

其中 \(i_{k}\) 是 13.3.2 中定义的同构。

例如，取 X 为 \(K^{n}\) 的一个开子集，并设 \(\{e_{1},\ldots,e_{n}\}\) 是 \(K^{n}\) 的典范基（与切空间 \(T_{x}(X)\) 恒同）。若 \(\alpha\in N^{n}\) 满足 \(|\alpha|\leqslant k\)，则算子 \(\Delta^{\alpha}\) 的符号由下列公式给出（参见 13.3.3）：

\[
\begin{array}{l l} \sigma_ {k} (\Delta^ {\alpha}) (x) = 0 & \text { if } \quad | \alpha | <   k \\ \sigma_ {k} (\Delta^ {\alpha}) (x) = \gamma_ {\alpha_ {1}} (e _ {1}) \ldots \gamma_ {\alpha_ {n}} (e _ {n}) & \text { if } \quad | \alpha | = k, \end{array}
\]

其中诸 \(\gamma_{\alpha_i}(e_i)\) 的乘积是对称积（参见 13.2.6 与 13.3.3）。

14.2.5（“复合的符号”）。记号与假设即 14.1.8 的那些。线性映射的复合定义了一个配对

\[
\mathcal {L} (\mathrm{F}; \mathrm{G}) \times \mathcal {L} (\mathrm{E}; \mathrm{F}) \rightarrow \mathcal {L} (\mathrm{E}; \mathrm{G}).
\]

另一方面，对称积运算定义了一个配对

\[
\mathsf {T S} ^ {k ^ {\prime \prime}} (\mathbf {T} (\mathbf {X})) \times \mathsf {T S} ^ {k ^ {\prime}} (\mathbf {T} (\mathbf {X})) \rightarrow \mathsf {T S} ^ {k} (\mathbf {T} (\mathbf {X})).
\]

由于

\[
\left. \begin{array}{l} \mathrm{S} ^ {k ^ {\prime \prime}} (\mathrm{F}, \mathrm{G}) = \mathrm{TS} ^ {k ^ {\prime \prime}} (\mathrm{T} (\mathrm{X})) \otimes \mathscr {L} (\mathrm{F}; \mathrm{G}) \\ \mathrm{S} ^ {k ^ {\prime}} (\mathrm{E}, \mathrm{F}) = \mathrm{TS} ^ {k ^ {\prime}} (\mathrm{T} (\mathrm{X})) \otimes \mathscr {L} (\mathrm{E}; \mathrm{F}) \\ \mathrm{S} ^ {k} (\mathrm{E}, \mathrm{G}) = \mathrm{TS} ^ {k} (\mathrm{T} (\mathrm{X})) \otimes \mathscr {L} (\mathrm{E}; \mathrm{G}) \end{array} \right\} \quad \text { cf. } 1 4. 2. 2,
\]

由此通过张量积得到配对

\[
\mathrm{S} ^ {k ^ {\prime \prime}} (\mathrm{F}, \mathrm{G}) \times \mathrm{S} ^ {k ^ {\prime}} (\mathrm{E}, \mathrm{F}) \rightarrow \mathrm{S} ^ {k} (\mathrm{E}, \mathrm{G}).\tag{*}
\]

于是有公式

\[
\sigma_ {k} (\mathbf {D} ^ {\prime \prime} \circ \mathbf {D} ^ {\prime}) = \sigma_ {k ^ {\prime \prime}} (\mathbf {D} ^ {\prime \prime}). \sigma_ {k ^ {\prime}} (\mathbf {D} ^ {\prime}),
\]

其中右端项中出现的乘积由上面的配对 (*) 定义。

现在假设 K 的特征为零，并设 \(x \in X\)。用 \(\sigma\)（相应地，\(\sigma'\)、\(\sigma''\)）表示 k-符号（相应地，\(k'\)-符号、\(k''\)-符号），这里指的是 \(D'' \circ D'\)（相应地，\(D'\)、\(D''\)）在 x 处的符号。于是，对每个 \(\omega \in \mathbf{T}_{x}(\mathbf{X})^{*}\)，有（沿用 14.2.2 末尾的记号）：

\[
\tilde {\sigma} (\omega) \in \mathscr {L} (\mathrm{E} _ {x}; \mathrm{G} _ {x}), \quad \tilde {\sigma} ^ {\prime} (\omega) \in \mathscr {L} (\mathrm{E} _ {x}; \mathrm{F} _ {x}), \quad \tilde {\sigma} ^ {\prime \prime} (\omega) \in \mathscr {L} (\mathrm{F} _ {x}; \mathrm{G} _ {x})
\]

以及

\[
\tilde {\sigma} (\omega) = \tilde {\sigma} ^ {\prime \prime} (\omega) \circ \tilde {\sigma} ^ {\prime} (\omega).
\]

